% ============================================================
\section{UML \& Vererbung (Abschnitt IV)}
% ============================================================
\textbf{UML} (Unified Modeling Language, aktuell 2.5.1, ISO/IEC 19505) ist eine grafische Modellierungssprache zur Konstruktion und Dokumentation von Software. Sie erg\"anzt Flussdiagramme: W\"ahrend Flussdiagramme \emph{sequentielle Abl\"aufe} zeigen, beantworten Klassendiagramme \emph{,,Wer macht was und wie h\"angt alles zusammen?''}.

\subsection{Klassendiagramme}
Eine Klasse ist ein Rechteck mit bis zu drei Abschnitten: \textbf{Name} / \textbf{Attribute} / \textbf{Methoden (Operationen)}. Eine zeitliche Abfolge wird nicht dargestellt.

\subsubsection*{Sichtbarkeit (Visibility)}
\begin{center}\small
\begin{tabular}{@{}cll@{}}
\toprule
\textbf{Symbol} & \textbf{Name} & \textbf{Sichtbar} \\
\midrule
\code{+} & public    & \"uberall \\
\code{-} & private   & nur in der Klasse \\
\code{\#} & protected & in Klasse \& abgeleiteten Klassen \\
\code{\textasciitilde} & package & im gesamten Paket \\
\bottomrule
\end{tabular}
\end{center}
Weitere Notation: Datentyp/R\"uckgabetyp nach \code{:}\,; \textbf{statische} Member werden \underline{unterstrichen}; \textbf{abgeleitete} Attribute (berechnet) mit \code{/} (z.\,B. \code{/alter}); erweiterte Eigenschaften in \code{\{ \}} (z.\,B. \code{readOnly}, \code{ordered}). Die Methodeneigenschaft \code{\{query\}} entspricht in C\texttt{++} einer \code{const}-Methode.

\begin{syntaxbox}[: Attribut / Methode in UML]
\code{sichtbarkeit name : Typ [= Default] \{Eigenschaften\}} \\
\code{+ berechne(in x : int) : double \{query\}}
\end{syntaxbox}

\subsection{Objektdiagramme}
Ein Objektdiagramm ist ein \emph{Snapshot} -- es zeigt konkrete Instanzen (Objekte) zu einem Zeitpunkt. Unterschiede zum Klassendiagramm: Objektname ist \underline{unterstrichen} (\code{myCar : Car}), Attribute haben nur \textbf{konkrete Werte}, Methoden werden \textbf{nicht} aufgelistet, keine Sichtbarkeit/Multiplizit\"at. Instanzen von Assoziationen hei\ss en \emph{Links}.

\subsection{Beziehungen zwischen Klassen}
Neben individuellen Eigenschaften definieren sich Objekte \"uber ihre \emph{Beziehungen}. Die drei Formen unterscheiden sich in der \textbf{St\"arke der Bindung}.

% --- Screenshot: Vergleich Assoziation/Aggregation/Komposition ---
\folie{0.85}{uml_beziehungen}{Assoziation, Aggregation und Komposition im Vergleich (Vorlesungsfolien Abschnitt IV).}

\begin{description}
  \item[Assoziation] allgemeine Beziehung (,,kennt'', ,,verwendet''). Lose gekoppelt, oft \"uber Pointer/Referenzen, \textbf{unabh\"angige} Lebensdauer. UML: einfache Linie, Pfeil zur bekannten Klasse.
  \item[Aggregation] ,,Ganzes-Teile''-Beziehung mit \textbf{schwachem} (nicht exklusivem) Besitz; Teile k\"onnen unabh\"angig existieren. UML: Linie mit \textbf{leerer Raute} auf der Seite des Ganzen.
  \item[Komposition] ,,besteht aus'' -- \textbf{starker, exklusiver} Besitz; Teile existieren nur mit dem Ganzen (gleiche Lebensdauer), meist direkte Member. UML: Linie mit \textbf{ausgef\"ullter Raute} auf der Seite des Ganzen.
\end{description}

\begin{lstlisting}
class Plane { Pilot* pilot; };          // Assoziation: kennt, unabhaengig
class Team  { Player* players; };       // Aggregation: hat, unabhaengig
class Plane { Engine engines[4]; };     // Komposition: besteht aus, gekoppelt
\end{lstlisting}

\subsubsection*{Multiplizit\"at}
Wie viele Instanzen stehen in Beziehung? Default ist \code{1}.
\begin{center}\small
\begin{tabular}{@{}ll@{}}
\toprule
\textbf{Notation} & \textbf{Bedeutung} \\
\midrule
\code{1} / \code{m} & genau eins / genau $m$ \\
\code{*} (\code{0..*}) & beliebig viele \\
\code{1..*} & mindestens eins \\
\code{0..1} & optional (0 oder 1) \\
\code{n..m} & Bereich (z.\,B. 2 bis 5) \\
\bottomrule
\end{tabular}
\end{center}

\subsection{Vererbung}
Vererbung modelliert eine \textbf{,,ist-ein''-Beziehung}: Eine \emph{Kind-/Unterklasse} (Derived) erbt Attribute und Methoden der \emph{Eltern-/Basisklasse} (Base) und kann sie erweitern oder spezialisieren.
\begin{syntaxbox}
\code{class Derived : public Base \{ ... \};}
\end{syntaxbox}

\begin{merke}
Der Zugriffsspezifizierer \code{protected} ist die Mischform: \code{protected}-Member sind in der eigenen Klasse \textbf{und} in abgeleiteten Klassen sichtbar -- aber nicht von au\ss en. \code{private}-Member sind \textbf{nie} in der abgeleiteten Klasse sichtbar.
\end{merke}

\subsubsection*{Vererbungsart (Zugriffsspezifizierer nach dem \texttt{:})}
Steuert, wie geerbte Member in der abgeleiteten Klasse sichtbar werden:
\begin{center}\small
\begin{tabular}{@{}lllll@{}}
\toprule
\textbf{Vererbung} & \code{public} aus Base & \code{protected} aus Base & \code{private} aus Base & \textbf{Bedeutung} \\
\midrule
\code{: public}    & bleibt public    & bleibt protected & nicht sichtbar & ,,Derived \emph{ist ein} Base'' \\
\code{: protected} & wird protected   & bleibt protected & nicht sichtbar & basiert auf Base (verborgen) \\
\code{: private}   & wird private     & wird private     & nicht sichtbar & nutzt Base intern \\
\bottomrule
\end{tabular}
\end{center}
\code{public} ist der Standard- und Normalfall.

\subsubsection*{Konstruktoren \& Destruktoren bei Vererbung}
\begin{itemize}
  \item Der Konstruktor der \textbf{Basisklasse wird zuerst} aufgerufen, dann der der abgeleiteten Klasse.
  \item Hat die Basisklasse \textbf{keinen} Default-Konstruktor, muss die abgeleitete Klasse explizit einen passenden Basiskonstruktor aufrufen.
  \item Destruktoren werden \textbf{nicht vererbt} und laufen in \textbf{umgekehrter} Reihenfolge (erst Derived, dann Base).
\end{itemize}
\begin{lstlisting}
class Base {
    int a;
public:
    Base(int a) : a(a) {}
};
class Derived : public Base {
public:
    Derived(int a) : Base(a) {}   // Basiskonstruktor explizit aufrufen (Pflicht!)
};
\end{lstlisting}
\begin{merke}
Reihenfolge: \textbf{Konstruktion} Base $\rightarrow$ Derived, \textbf{Destruktion} Derived $\rightarrow$ Base. In UML hei\ss t Vererbung \emph{Generalisierung}: Linie mit \textbf{leerer Dreieckspitze}, die stets auf die \textbf{Basisklasse} zeigt.
\end{merke}
\clearpage
